\documentclass[11pt,a4paper]{article}
\usepackage{fontspec}
\usepackage{xeCJK}
\setCJKmainfont{STSong}
\setCJKsansfont{STHeiti}
\usepackage{amsmath,amssymb,amsthm}
\usepackage{mathtools}
\usepackage{bm}
\usepackage{booktabs}
\usepackage{array}
\usepackage{enumitem}
\usepackage{geometry}
\usepackage{xcolor}
\usepackage{tcolorbox}
\usepackage{algorithm}
\usepackage{algpseudocode}
\usepackage{pifont}
\usepackage{tikz}
\usetikzlibrary{arrows.meta,positioning,calc,decorations.pathreplacing}
\usepackage[pdfencoding=auto,psdextra]{hyperref}

\geometry{margin=2.5cm}

\newcommand{\loss}{\mathcal{L}}
\newcommand{\E}{\mathbb{E}}
\newcommand{\R}{\mathbb{R}}
\newcommand{\KL}{\mathrm{KL}}
\newcommand{\N}{\mathcal{N}}
\newcommand{\adv}{\mathcal{A}}
\newcommand{\reward}{R}
\newcommand{\softmax}{\mathrm{softmax}}
\newcommand{\sg}{\mathrm{sg}}

\theoremstyle{definition}
\newtheorem{definition}{Definition}
\newtheorem{remark}{Remark}
\newtheorem{proposition}{Proposition}
\newtheorem{corollary}{Corollary}
\newtheorem{lemma}{Lemma}
\newtheorem{theorem}{Theorem}
\newtheorem{example}{Example}
\newtheorem{strategy}{Strategy}

\title{\textbf{MoTV Phase 2: Per-Prompt-Set $\lambda$ 架构分析}\\[0.3em]
\large $\ell \in \R^{K \times T \times S}$ 的 Loss、梯度与 Reward 设计}
\author{Flow Factory}
\date{}

\begin{document}
\maketitle

\begin{abstract}
基于 Phase~1 的分析结论，本文正式采用 per-prompt-set 扩展的 MoTV 架构：logits $\ell \in \R^{K \times T \times S}$，其中 $K$ = teacher 数量，$T$ = 推理 timesteps，$S$ = prompt set 数量。每个 prompt set 拥有独立的 teacher 混合策略。本文给出完整的 loss 推导、梯度公式、收敛性分析，并详细设计 per-set reward 策略——包括 reward normalization、advantage computation、OOD bonus 的精确公式和实现细节。
\end{abstract}

\tableofcontents
\newpage

%% ============================================================
\section{架构定义}
\label{sec:architecture}
%% ============================================================

\subsection{符号表}

\begin{center}
\renewcommand{\arraystretch}{1.3}
\begin{tabular}{cl}
\toprule
\textbf{符号} & \textbf{含义} \\
\midrule
$K$ & Teacher 数量（= 3: GenEval, PickScore, OCR） \\
$T$ & 推理 timesteps 数量（= \texttt{num\_inference\_steps}） \\
$S$ & Prompt set 数量（= 3: $\mathcal{P}_{\mathrm{gen}}, \mathcal{P}_{\mathrm{pick}}, \mathcal{P}_{\mathrm{ocr}}$） \\
$B$ & Batch size \\
$d$ & Latent 维度 $= C \times H \times W$ \\
$\ell \in \R^{K \times T \times S}$ & 可学习 logits（总参数量 $= K \times T \times S$） \\
$\tau$ & Softmax temperature \\
$\lambda_{k,i,s}$ & Teacher $k$ 在 timestep $i$、prompt set $s$ 上的权重 \\
$v^{(k)}(x, t)$ & Teacher $k$ 的 frozen velocity field \\
$v_{\lambda_s}(x, t_i)$ & Prompt set $s$ 上的 combined velocity \\
$R_j(x)$ & 第 $j$ 个 reward function（GenEval / PickScore / OCR） \\
$\adv_s$ & Prompt set $s$ 上的 advantage \\
$\beta$ & NFT beta（positive/negative interpolation） \\
$\gamma$ & OOD bonus 系数 \\
\bottomrule
\end{tabular}
\end{center}

\subsection{Per-Set Velocity Field}

\begin{definition}[Per-Set Combined Velocity]
对 prompt $p \in \mathcal{P}_s$（属于 prompt set $s$），定义：
\begin{equation}
\boxed{
v_{\lambda_s}(x, t_i) = \sum_{k=1}^K \lambda_{k,i,s}\, v^{(k)}(x, t_i), \quad
\lambda_{k,i,s} = \frac{\exp(\ell_{k,i,s}/\tau)}{\sum_{k'=1}^K \exp(\ell_{k',i,s}/\tau)}
}
\label{eq:per_set_velocity}
\end{equation}
其中权重满足 $\lambda_{k,i,s} \geq 0$, $\sum_{k=1}^K \lambda_{k,i,s} = 1$, $\forall\, i, s$。
\end{definition}

\begin{remark}[与 Phase~1 Global $\lambda$ 的关系]
Global $\lambda$ 是 per-set $\lambda$ 的特例：当 $\ell_{k,i,s} = \ell_{k,i}$ 对所有 $s$ 相同时，退化为 Phase~1 的 global 版本。Per-set 扩展引入了 $S$ 倍的参数，但仍然极小（$K \times T \times S = 3 \times 28 \times 3 = 252$ 参数）。
\end{remark}

\subsection{参数空间几何}

\begin{proposition}[搜索空间结构]
Per-set MoTV 的搜索空间为 $S$ 个独立的乘积 simplex：
\begin{equation}
\Omega = \prod_{s=1}^S \prod_{i=1}^T \Delta^K = \left(\Delta^K\right)^{T \times S}
\end{equation}
其中 $\Delta^K = \{\lambda \in \R^K_+ : \sum_k \lambda_k = 1\}$ 是 $K$-simplex。

\textbf{关键性质}：不同 prompt set $s$ 和 $s'$ 的参数 $\ell_{:,:,s}$ 和 $\ell_{:,:,s'}$ 完全独立——梯度不交叉，优化完全解耦。
\end{proposition}

\subsection{Inference 流程}

\begin{algorithm}[H]
\caption{Per-Set MoTV Inference}
\label{alg:inference}
\begin{algorithmic}[1]
\Require Prompt $p$, noise $x_T \sim \N(0, I)$, logits $\ell \in \R^{K \times T \times S}$
\Ensure Generated image $x_0$
\State Determine prompt set: $s \gets \mathrm{classify}(p)$
\State Compute weights: $\lambda_{:,:,s} \gets \softmax_k(\ell_{:,:,s} / \tau)$ \Comment{shape $(K, T)$}
\State $x \gets x_T$
\For{$i = 1, \ldots, T$}
    \For{$k = 1, \ldots, K$}
        \State $v^{(k)}_i \gets v^{(k)}(x, t_i)$ \Comment{Frozen teacher forward}
    \EndFor
    \State $v_{\mathrm{combined}} \gets \sum_k \lambda_{k,i,s} \cdot v^{(k)}_i$ \Comment{Weighted combination}
    \State $x \gets \mathrm{scheduler\_step}(x, v_{\mathrm{combined}}, t_i)$ \Comment{ODE step}
\EndFor
\State \Return $x$
\end{algorithmic}
\end{algorithm}

%% ============================================================
\section{Loss 函数}
\label{sec:loss}
%% ============================================================

\subsection{Per-Set DiffusionNFT Loss}

对每个 prompt set $s$ 独立定义 NFT loss。设 batch 中 prompt $p \in \mathcal{P}_s$：

\begin{definition}[Per-Set NFT Loss (single timestep)]
\begin{equation}
\boxed{
\loss_{i,s}(\ell_{:,:,s}) = \frac{1}{\beta}\left[
r_s \cdot L^+_{i,s} + (1 - r_s) \cdot L^-_{i,s}
\right]
}
\label{eq:per_set_nft_loss}
\end{equation}
其中：
\begin{align}
r_s &= \frac{\adv_s + A_{\max}}{2 A_{\max}} \in [0, 1] \label{eq:normalized_adv} \\
v^{\mathrm{new}}_{i,s} &= \sum_k \lambda_{k,i,s}\, v^{(k)}(x_{t_i}, t_i) \quad\text{(current policy for set $s$)} \label{eq:v_new_s} \\
v^{\mathrm{old}}_{i,s} &= \sum_k \bar{\lambda}_{k,i,s}\, v^{(k)}(x_{t_i}, t_i) \quad\text{(EMA policy for set $s$)} \label{eq:v_old_s} \\
v^+_{i,s} &= \beta\, v^{\mathrm{new}}_{i,s} + (1-\beta)\, v^{\mathrm{old}}_{i,s} \label{eq:v_plus} \\
v^-_{i,s} &= (1+\beta)\, v^{\mathrm{old}}_{i,s} - \beta\, v^{\mathrm{new}}_{i,s} \label{eq:v_minus}
\end{align}
\begin{align}
x^+_{0,s} &= x_{t_i} - \sigma_i\, v^+_{i,s} \label{eq:x0_plus} \\
x^-_{0,s} &= x_{t_i} - \sigma_i\, v^-_{i,s} \label{eq:x0_minus} \\
L^+_{i,s} &= \frac{\|x^+_{0,s} - x_0\|^2}{(w^+_s)^2}, \quad
L^-_{i,s} = \frac{\|x^-_{0,s} - x_0\|^2}{(w^-_s)^2} \label{eq:L_plus_minus}
\end{align}
$w^+_s, w^-_s$ 为 detached mean absolute error（per-sample normalization）。
\end{definition}

\subsection{Full Training Loss}

累积所有 timestep 和 batch samples：
\begin{equation}
\loss_s(\ell_{:,:,s}) = \frac{1}{B_s \cdot T}\sum_{b \in \mathcal{B}_s}\sum_{i=1}^T \loss^{(b)}_{i,s}(\ell_{:,:,s})
\label{eq:full_loss}
\end{equation}
其中 $\mathcal{B}_s$ 为当前 batch 中属于 prompt set $s$ 的 sample 子集，$B_s = |\mathcal{B}_s|$。

\begin{remark}[Per-Set 独立性]
总 loss 为各 set 的 loss 之和：
\begin{equation}
\loss_{\mathrm{total}}(\ell) = \sum_{s=1}^S \loss_s(\ell_{:,:,s})
\end{equation}
由于 $\ell_{:,:,s}$ 只出现在 $\loss_s$ 中，$\frac{\partial \loss_{\mathrm{total}}}{\partial \ell_{:,:,s}} = \frac{\partial \loss_s}{\partial \ell_{:,:,s}}$。各 prompt set 的优化完全独立——如同并行训练 $S$ 个独立的 global-$\lambda$ MoTV。
\end{remark}

%% ============================================================
\section{梯度推导}
\label{sec:gradient}
%% ============================================================

\subsection{Main Theorem}

\begin{theorem}[Per-Set MoTV Gradient]
\label{thm:per_set_gradient}
对 prompt set $s$，single sample，single timestep $i$：
\begin{equation}
\boxed{
\frac{\partial \loss_{i,s}}{\partial \ell_{k,i,s}} =
\frac{2}{\tau}\,\lambda_{k,i,s}\left[v^{(k)}_i - v^{\mathrm{new}}_{i,s}\right]^\top
\underbrace{\left[\frac{r_s\,(x^+_{0,s} - x_0)}{(w^+_s)^2} - \frac{(1-r_s)\,(x^-_{0,s} - x_0)}{(w^-_s)^2}\right]}_{\triangleq\; \mathbf{e}_{i,s} \;\in\; \R^d}\cdot(-\sigma_i)
}
\label{eq:per_set_gradient}
\end{equation}

且 $\frac{\partial \loss_{i,s}}{\partial \ell_{k,j,s'}} = 0$ 当 $j \neq i$ 或 $s' \neq s$。
\end{theorem}

\begin{proof}
推导与 Phase~1 Theorem~2.1 完全相同，仅将 global $\lambda_{k,i}$ 替换为 per-set $\lambda_{k,i,s}$。各 set 的 sample 只对自己的 $\ell_{:,:,s}$ 产生梯度，跨 set 梯度为零。
\end{proof}

\subsection{三重解耦性}

\begin{corollary}[Triple Decoupling]
\label{cor:triple_decouple}
Per-set MoTV 的梯度具有三重解耦：
\begin{enumerate}[leftmargin=1.5em]
\item \textbf{Timestep 解耦}：$\frac{\partial \loss_{i,s}}{\partial \ell_{k,j,s}} = 0$ for $j \neq i$
\item \textbf{Prompt-set 解耦}：$\frac{\partial \loss_{i,s}}{\partial \ell_{k,i,s'}} = 0$ for $s' \neq s$
\item \textbf{Orthogonal teacher competition}：在每个 $(i, s)$ 固定的 slice 上，$K$ 个 teacher 通过 softmax 竞争。增大 $\ell_{k,i,s}$ 必然减小 $\ell_{k',i,s}$ ($k' \neq k$) 的有效权重。
\end{enumerate}

\textbf{含义}：每个 $(\text{timestep}, \text{prompt set})$ pair 是一个完全独立的 $K$-way 选择问题。优化 252 个参数等价于并行优化 $T \times S = 84$ 个独立的 $K$-simplex 选择。
\end{corollary}

\subsection{Gradient 的分量分析}

将梯度~\eqref{eq:per_set_gradient} 分解为三个语义明确的因子：

\begin{equation}
\frac{\partial \loss_{i,s}}{\partial \ell_{k,i,s}} = \underbrace{\frac{2\sigma_i}{\tau}\,\lambda_{k,i,s}}_{\text{(I) Scale}} \;\cdot\;
\underbrace{\left[v^{(k)}_i - v^{\mathrm{new}}_{i,s}\right]}_{\text{(II) Direction}} \;\cdot\;
\underbrace{(-\mathbf{e}_{i,s})}_{\text{(III) Reward signal}}
\label{eq:grad_decomp}
\end{equation}

\paragraph{Factor (I): Scale.}
\begin{itemize}[leftmargin=1.5em]
\item $\lambda_{k,i,s} \in [0, 1]$：当前权重。权重越接近 0 或 1，梯度越小（softmax 饱和效应）。最大梯度在 $\lambda = 1/K$ 处（均匀分布时，各 teacher 梯度最灵敏）。
\item $\sigma_i$：noise level at timestep $i$。高噪声区域（早期 timestep，$\sigma \approx 1$）梯度较大；低噪声区域（晚期，$\sigma \to 0$）梯度较小——这意味着\textbf{早期 timestep 的权重学习更快}。
\end{itemize}

\paragraph{Factor (II): Direction $[v^{(k)}_i - v^{\mathrm{new}}_{i,s}]$.}
\begin{itemize}[leftmargin=1.5em]
\item Teacher $k$ 相对于当前 combined velocity 的``偏差方向''。
\item 若 $\lambda_{k,i,s} = 1$（已选中 teacher $k$），则 $v^{\mathrm{new}}_{i,s} = v^{(k)}_i$，此项为 0——\textbf{已确定的选择不再变化}。
\item 若 $\lambda_{k,i,s} \ll 1$，则 $v^{(k)}_i - v^{\mathrm{new}}_{i,s} \approx v^{(k)}_i - \sum_{k'\neq k}\lambda_{k'} v^{(k')}_i$——teacher $k$ 与其他 teacher 混合的差异。
\end{itemize}

\paragraph{Factor (III): Reward signal $\mathbf{e}_{i,s}$.}
\begin{itemize}[leftmargin=1.5em]
\item 由 advantage $r_s$ 加权的 prediction error。
\item 当 $r_s > 0.5$（正 advantage，当前 sample 好）：$\mathbf{e}_{i,s}$ 主要由 $L^+$ 方向决定。
\item 当 $r_s < 0.5$（负 advantage，当前 sample 差）：$\mathbf{e}_{i,s}$ 主要由 $-L^-$ 方向决定。
\end{itemize}

\subsection{简化梯度（$\beta = 1$, $w = 1$, $v^{\mathrm{old}} \approx v^{\mathrm{new}}$）}

在标准设定下（$\beta = 1$, 忽略 normalization weight，small off-policy gap）：
\begin{equation}
\boxed{
\frac{\partial \loss_{i,s}}{\partial \ell_{k,i,s}} \approx \frac{2\sigma_i^2}{\tau}\,\lambda_{k,i,s}
\left[v^{(k)}_i - v^{\mathrm{new}}_{i,s}\right]^\top
\left[(2r_s - 1)(v^{\mathrm{new}}_{i,s} - v^*_i)\right]
}
\label{eq:simplified_grad}
\end{equation}
其中 $v^*_i = (x_{t_i} - x_0)/\sigma_i$ 为 ground truth velocity。

\begin{tcolorbox}[colback=green!5, colframe=green!50, title=梯度动力学解读]
设 $\adv_s > 0$（好 sample，$r_s > 0.5$）：
\begin{itemize}[leftmargin=1em]
\item 若 teacher $k$ 的 velocity 比当前混合更接近 ground truth（$[v^{(k)} - v^{\mathrm{new}}]^\top(v^{\mathrm{new}} - v^*) < 0$），则梯度为负 → gradient descent 增大 $\ell_{k,i,s}$ → \textbf{增加 teacher $k$ 在 set $s$, timestep $i$ 上的权重}。
\item 若 teacher $k$ 比混合更差，则权重减小。
\end{itemize}

设 $\adv_s < 0$（差 sample，$r_s < 0.5$）：
\begin{itemize}[leftmargin=1em]
\item 反向信号：惩罚当前混合方向，探索远离当前配置的方向。
\end{itemize}

\textbf{Per-set 的额外优势}：set $s$ 的梯度只看 set $s$ 的 reward 信号——GenEval set 上的 reward 只影响 $\ell_{:,:,\mathrm{gen}}$，OCR set 上的 reward 只影响 $\ell_{:,:,\mathrm{ocr}}$。不同 set 的 reward 信号\textbf{完全互不干扰}。
\end{tcolorbox}

%% ============================================================
\section{收敛性与最优解分析}
\label{sec:convergence}
%% ============================================================

\subsection{Per-Set 最优解}

\begin{theorem}[Per-Set 最优解的形式]
\label{thm:optimal}
对 prompt set $s$，optimal logits $\ell^*_{:,:,s}$ 满足：
\begin{equation}
\ell^*_{:,:,s} = \arg\max_{\ell_{:,:,s}} \E_{p \in \mathcal{P}_s}\left[\reward_s\!\left(\mathrm{ODE}(v_{\lambda_s}; x_T, p)\right)\right]
\label{eq:per_set_optimal}
\end{equation}
其中 $\reward_s$ 是 prompt set $s$ 的 reward function（定义见 Section~\ref{sec:reward}）。
\end{theorem}

\subsection{Level~1 保证（$\geq$ Single Teacher）}

\begin{theorem}[Guaranteed Non-Degradation]
\label{thm:non_degrade}
对任意 prompt set $s$ 和 teacher $k$：
\begin{equation}
\max_{\ell_{:,:,s}} \E_{p \in \mathcal{P}_s}\left[\reward_s(v_{\lambda_s})\right] \geq \reward_s(v^{(k)}), \quad \forall k = 1,\ldots,K
\end{equation}

特别地，取 $k = s$（in-domain teacher）：
\begin{equation}
\boxed{
\E[\reward_s(v_{\lambda^*_s})] \geq \E[\reward_s(v^{(s)})] \quad\text{(in-domain performance $\geq$ single teacher)}
}
\end{equation}
\end{theorem}

\begin{proof}
$\lambda_{k,i,s} = \delta_{k,s}$（teacher $s$ 恒取权重 1）是可行解，对应 $v_{\lambda_s} = v^{(s)}$。优化在包含此点的搜索空间上进行，最优解不劣于此。对任意其他 teacher $k$ 同理（$\delta_{k,\cdot}$ 也是可行解）。
\end{proof}

\begin{remark}[与 OPD 的对比]
这是 MoTV per-set 架构相比 OPD-Average 的根本优势。OPD-Average 的最优解是 teacher 均值，由 Jensen 不等式必然在 in-domain 上退化。Per-set MoTV 则有\textbf{构造性的下界保证}。
\end{remark}

\subsection{Level~2 条件（$>$ Single Teacher）}

\begin{theorem}[Strict Improvement Condition]
\label{thm:strict_improve}
Per-set MoTV 严格超越 teacher $s$ on set $s$ 的充分条件为：存在 timestep $i^*$ 和 teacher $k^* \neq s$ 使得：
\begin{equation}
\E_{p \in \mathcal{P}_s}\left[\reward_s\!\left(\mathrm{ODE}(v_{\lambda'}; x_T, p)\right)\right] > \E_{p \in \mathcal{P}_s}\left[\reward_s\!\left(\mathrm{ODE}(v^{(s)}; x_T, p)\right)\right]
\end{equation}
其中 $\lambda'$ 仅在 timestep $i^*$ 处与 $\delta_{k,s}$ 不同：$\lambda'_{k^*,i^*,s} = \epsilon > 0$, $\lambda'_{s,i^*,s} = 1 - \epsilon$。

\textbf{直觉}：只需存在一个 timestep，在该 timestep 上引入少量其他 teacher 的 velocity 能提升 set $s$ 的 reward。
\end{theorem}

\subsection{温度 $\tau$ 与收敛形态}

\begin{proposition}[温度对收敛的影响]
\label{prop:temperature}
\begin{enumerate}[leftmargin=1.5em]
\item $\tau \to 0$（低温极限）：$\lambda_{k,i,s} \to \delta_{k, \arg\max_k \ell_{k,i,s}}$。每个 $(i,s)$ 退化为 hard selection（winner-take-all）。搜索空间退化为 $\{1,\ldots,K\}^{T \times S}$——离散组合优化。

\item $\tau \to \infty$（高温极限）：$\lambda_{k,i,s} \to 1/K$。所有 teacher 均匀混合，等价于 teacher velocity 的算术平均。

\item $\tau$ 适中（推荐 $\tau \in [0.5, 2.0]$）：soft assignment，允许梯度流动的同时保持一定的 specialization tendency。
\end{enumerate}
\end{proposition}

\begin{remark}[温度退火策略]
实践建议：初始 $\tau = 1.0$（充分探索），训练后期逐步降低到 $\tau = 0.3$--$0.5$（promote specialization）。这类似 Gumbel-Softmax 的退火策略。
\end{remark}

%% ============================================================
\section{Reward 设计（详细版）}
\label{sec:reward}
%% ============================================================

\subsection{Per-Set Composite Reward Definition}

\begin{definition}[Per-Set Composite Reward with OOD Bonus]
\label{def:per_set_reward}
对 prompt $p \in \mathcal{P}_s$ 生成的 image $x$，设 $\mathcal{R}(x, s)$ 为该 sample 的 applicable reward 集合（由 \texttt{applicable\_sources} 过滤，不适用的 reward 得到 NaN 被跳过）：
\begin{equation}
\boxed{
r_{\mathrm{composite}}(x, s) = R_{\mathrm{in}(s)}(x) + \gamma \cdot \frac{1}{|\mathcal{R}_{\mathrm{ood}}|}\sum_{j \in \mathcal{R}_{\mathrm{ood}}} R_j(x)
}
\label{eq:reward_def}
\end{equation}
其中 $\mathrm{in}(s)$ 是 source $s$ 的 in-domain reward（由 \texttt{TeacherConfig.reward\_name} 显式指定），$\mathcal{R}_{\mathrm{ood}} = \mathcal{R}(x,s) \setminus \{\mathrm{in}(s)\}$，$\gamma \geq 0$ 是 OOD bonus 系数。

注意：这里使用的是 \textbf{raw reward}（未做 normalization）。Reward scale 的差异通过下游 advantage 计算中的 group-std normalization 隐式消除。
\end{definition}

\begin{remark}[Reward 权重解读]
\begin{itemize}[leftmargin=1.5em]
\item In-domain reward $R_{\mathrm{in}(s)}$ 的有效系数为 $1$。
\item 每个 OOD reward $R_j$ ($j \neq \mathrm{in}(s)$) 的有效系数为 $\frac{\gamma}{|\mathcal{R}_{\mathrm{ood}}|}$。
\item 当 $\gamma = 0$：纯 in-domain 优化（Level~1 保证最强，但可能收敛到 teacher $s$ 本身）。
\item 当 $\gamma = 0.2$（推荐）：OOD 信号提供探索方向但不干扰主信号。
\end{itemize}
\end{remark}

\subsection{Advantage Computation (GDPO-Style Per-Reward Normalization)}

MoF 采用 GDPO 的核心思想：\textbf{先对每个 reward 独立做 group-level normalization 得到 per-reward advantage，然后在 advantage 空间做 in-domain + OOD 加权聚合}。

\begin{definition}[MoF Advantage Pipeline]
\label{def:advantage}
设 epoch 中有 $N$ 个 sample（跨所有 rank gather 后），按 prompt 分组为 $\mathcal{G}_1, \ldots, \mathcal{G}_M$（每个 group 有 $K$ 个 sample）。

\textbf{Step 1: Per-Reward Group Normalization.} 对每个 reward $k$，独立计算其 advantage：
\begin{equation}
a_k(x_i) = \frac{R_k(x_i) - \mu_{k,g}}{\max(\sigma_{k,g},\, \epsilon)}, \quad
\mu_{k,g} = \frac{1}{|\mathcal{G}_g|}\sum_{j \in \mathcal{G}_g} R_k(x_j), \quad
\sigma_{k,g} = \mathrm{std}(\{R_k(x_j)\}_{j \in \mathcal{G}_g})
\label{eq:per_reward_norm}
\end{equation}
此步将各 reward 统一到 $\sim$ zero-mean, unit-variance 的 advantage 空间，\textbf{消除了 reward scale 差异}（如 GenEval $\in [0,1]$ vs PickScore $\in [0.8, 0.95]$）。

\textbf{Step 2: Per-Set Advantage Aggregation.} 对 sample $x_i$ 来自 source $s$：
\begin{equation}
\boxed{
\mathcal{A}_{\mathrm{combined}}(x_i) = a_{\mathrm{in}(s)}(x_i) + \gamma \cdot \frac{1}{|\mathcal{R}_{\mathrm{ood}}|}\sum_{j \in \mathcal{R}_{\mathrm{ood}}} a_j(x_i)
}
\label{eq:advantage_aggregation}
\end{equation}
其中 $\mathcal{R}_{\mathrm{ood}} = \mathcal{R}(x_i, s) \setminus \{\mathrm{in}(s)\}$ 是适用于该 sample 的非 in-domain reward 集合。

\textbf{Step 3: Batch-Level Normalization.}
\begin{equation}
\mathcal{A}_{\mathrm{final}}(x_i) = \frac{\mathcal{A}_{\mathrm{combined}}(x_i)}{\max(\sigma_{\mathrm{global}},\, \epsilon)}, \quad
\sigma_{\mathrm{global}} = \mathrm{std}(\{\mathcal{A}_{\mathrm{combined}}(x_i)\}_{i=1}^N)
\label{eq:final_norm}
\end{equation}
\end{definition}

\begin{remark}[与原始 GDPO 的关系]
原始 GDPO 的公式为：
\begin{equation}
\mathcal{A}^{\mathrm{GDPO}}_i = \frac{\sum_k w_k \cdot a_k(x_i)}{\sigma_{\mathrm{global}}}
\end{equation}
其中 $w_k$ 是全局静态权重。MoF 的区别在于：\textbf{权重是 per-sample 自适应的}——对来自 source $s$ 的 sample，in-domain reward 权重为 $1$，OOD reward 权重为 $\gamma/|\mathcal{R}_{\mathrm{ood}}|$。这等价于 GDPO with sample-conditional weights。
\end{remark}

\begin{remark}[为什么在 Advantage 空间聚合而非 Raw Reward 空间]
如果在 raw reward 空间聚合（$r = R_{\mathrm{in}} + \gamma \cdot R_{\mathrm{ood}}$），然后做 group normalization，会有以下问题：
\begin{itemize}[leftmargin=1.5em]
\item GenEval reward 范围 $[0, 1]$，PickScore 范围 $[0.8, 0.95]$——scale 不同导致 $\gamma$ 的有效影响力不可预测。
\item 同一 $\gamma = 0.2$ 对不同 reward pair 的实际效果差异巨大。
\end{itemize}
在 advantage 空间聚合（各 reward 已归一化为 $\sim N(0, 1)$）后：
\begin{itemize}[leftmargin=1.5em]
\item $\gamma = 0.2$ 意味着 OOD advantage 贡献 20\% 的信号强度，\textbf{与 reward 原始 scale 无关}。
\item 不需要手动调整 reward weights 来补偿 scale 差异。
\end{itemize}
\end{remark}

\begin{remark}[Cross-Source Group 不混合的 Invariant]
关键约束：同一 group 内的 $K$ 个 sample 必须来自同一 source（因为共享同一 prompt）。因此 per-reward group normalization 中，每个 group 内所有 sample 使用相同的 composite 公式——不会出现 group 内有的 sample 有某个 reward 而有的没有的情况。NaN（非 applicable）只发生在\textbf{跨 group} 维度。
\end{remark}

\subsection{OOD Bonus $\gamma$ 的梯度效应}

\begin{proposition}[OOD Bonus 如何引导探索]
\label{prop:ood_bonus}
设当前 $\lambda^*_{:,:,s} \approx \delta_{k,s}$（已收敛到 in-domain teacher）。引入 $\gamma > 0$ 后，梯度为：
\begin{align}
\frac{\partial \loss_{i,s}}{\partial \ell_{k',i,s}} &\propto \lambda_{k',i,s}\left[v^{(k')}_i - v^{(s)}_i\right]^\top \mathbf{e}_{i,s}
\end{align}
其中 $\mathbf{e}_{i,s}$ 现在包含 OOD reward 的贡献。

当 OOD teacher $k'$ 在某个 timestep $i$ 上能提升 OOD reward $R_{k'}$ 而不显著损害 $R_{\mathrm{in}(s)}$ 时：
\begin{itemize}[leftmargin=1.5em]
\item $\mathbf{e}_{i,s}$ 的 OOD 分量驱动 $\ell_{k',i,s}$ 增大。
\item $\lambda_{k',i,s}$ 从 $\sim 0$ 缓慢增长，引入少量 teacher $k'$ 的 velocity。
\item 若这确实提升了综合 reward $r_{\mathrm{composite}}$，则正向 advantage 强化这一趋势。
\end{itemize}

\textbf{本质}：$\gamma > 0$ 为已收敛的 solution 提供了``突破 single-teacher 局部最优''的梯度推力。
\end{proposition}

%% ============================================================
\section{$\gamma$ 的理论选择}
\label{sec:gamma_theory}
%% ============================================================

\subsection{$\gamma$ 过大的风险}

\begin{proposition}[$\gamma$ Upper Bound]
\label{prop:gamma_upper}
若 $\gamma$ 过大（OOD reward 占主导），则最优 $\lambda^*_{:,:,s}$ 将偏离 in-domain teacher，导致 in-domain performance 退化。具体地，当：
\begin{equation}
\gamma > \frac{\Delta R_{\mathrm{in}}}{\Delta R_{\mathrm{ood}}}
\end{equation}
其中 $\Delta R_{\mathrm{in}}$ = 从 teacher $s$ 偏离时 $R_s$ 的下降量，$\Delta R_{\mathrm{ood}}$ = 同时 OOD reward 的提升量，则 in-domain 会退化。

\textbf{经验估计}：从实验数据看，OPD-Average（完全均匀混合）导致 in-domain 下降 $\sim 0.1$，OOD 提升 $\sim 0.01$--$0.03$。因此 $\Delta R_{\mathrm{in}}/\Delta R_{\mathrm{ood}} \sim 3$--$10$。安全的 $\gamma$ 应远小于此。
\end{proposition}

\subsection{推荐 $\gamma$ 范围}

\begin{tcolorbox}[colback=blue!5, colframe=blue!40, title=推荐配置]
\begin{center}
\begin{tabular}{lcc}
\toprule
\textbf{策略} & $\gamma$ & \textbf{预期效果} \\
\midrule
保守（Level~1 为主） & $0$ & 确保 $\geq$ single teacher，无 OOD bonus \\
推荐（Level~1 + Level~2 探索） & $0.2$ & in-domain 安全，OOD 有探索空间 \\
激进（Level~2 为主） & $0.5$ & 更强 OOD 探索，in-domain 可能微降 \\
\bottomrule
\end{tabular}
\end{center}

\textbf{建议}：先用 $\gamma = 0$ 验证 Level~1 保证（确认 MoTV 至少达到 single teacher），再切换 $\gamma = 0.2$ 探索 Level~2。
\end{tcolorbox}

%% ============================================================
\section{实现细节}
\label{sec:implementation}
%% ============================================================

\subsection{Prompt Set Classification}

\begin{strategy}[Prompt Set Determination]
\label{strat:classify}
每个 prompt 携带 metadata 标注其 source set。具体来说，dataloader 中每个 batch element 包含 field \texttt{prompt\_set\_id} $\in \{0, 1, \ldots, S-1\}$。

\textbf{实现}：在 \texttt{MoTVTrainer.optimize()} 中：
\begin{enumerate}[leftmargin=1.5em]
\item 从 batch 中提取 \texttt{prompt\_set\_id} tensor (shape $(B,)$)
\item 对每个 set $s$：$\mathcal{B}_s = \{b : \texttt{prompt\_set\_id}[b] = s\}$
\item 使用 $\ell_{:,:,s}$ 对 $\mathcal{B}_s$ 中的 sample 计算 loss 和梯度
\end{enumerate}
\end{strategy}

\subsection{Logits 初始化}

\begin{strategy}[Initialization for Per-Set Logits]
\label{strat:init}
两种初始化策略：

\textbf{Option A: Zeros}（推荐用于验证 Level~1）
\begin{equation}
\ell_{k,i,s} = 0, \quad \forall k, i, s
\end{equation}
初始权重为均匀 $\lambda_{k,i,s} = 1/K$。所有 teacher 起点平等。

\textbf{Option B: Teacher-Biased}（推荐用于 Level~2 探索）
\begin{equation}
\ell_{k,i,s} = \begin{cases}
C & \text{if } k = s \\
0 & \text{otherwise}
\end{cases}, \quad C \approx 2.0
\end{equation}
初始权重偏向 in-domain teacher：$\lambda_{s,i,s} \approx 0.78$（$K=3$, $C=2$），其他 teacher $\approx 0.11$。从 single teacher 出发，通过 reward-driven optimization 逐步引入其他 teacher 的贡献。
\end{strategy}

\begin{remark}[Option B 的动机]
Option B 的初始点即为``near single-teacher'' solution，确保训练初期不退化。$\gamma > 0$ 的 OOD bonus 从此基础上驱动探索。若 teacher $k' \neq s$ 在某些 timestep 上有正贡献，gradient 会逐步提升 $\ell_{k',i,s}$。
\end{remark}

\subsection{EMA 与 Off-Policy}

\begin{strategy}[Per-Set EMA]
对整个 logits tensor $\ell \in \R^{K \times T \times S}$ 维护 EMA：
\begin{equation}
\bar{\ell} \leftarrow (1 - \alpha_{\mathrm{ema}})\,\bar{\ell} + \alpha_{\mathrm{ema}}\,\ell
\end{equation}
其中 $\alpha_{\mathrm{ema}} = 1 - \texttt{logits\_ema\_decay}$（如 decay $= 0.99$ → $\alpha = 0.01$）。

Sampling 使用 $\bar{\ell}$（off-policy），training 使用 $\ell$（on-policy），产生 NFT 的 new/old policy gap。
\end{strategy}

\subsection{Training Batch 组织}

\begin{strategy}[Mixed-Set Batch]
\label{strat:batch}
每个 training batch 包含来自\textbf{所有} $S$ 个 prompt set 的 sample（按比例混合）：
\begin{equation}
B = B_{\mathrm{gen}} + B_{\mathrm{pick}} + B_{\mathrm{ocr}}, \quad B_s \approx B / S
\end{equation}

\textbf{Batch 内并行}：一次 forward pass 产生所有 sample 的 teacher velocities，然后按 \texttt{prompt\_set\_id} 分组计算各 set 的 loss 和梯度。

\textbf{梯度累积}：由于三个 set 的梯度作用于 $\ell$ 的不同 slice（$\ell_{:,:,0}$, $\ell_{:,:,1}$, $\ell_{:,:,2}$），可以直接 in-place 累积，无冲突。
\end{strategy}

\subsection{Optimizer 配置}

\begin{strategy}[Per-Set Optimizer]
虽然参数解耦，但实践中使用\textbf{单个 AdamW optimizer} 管理整个 $\ell \in \R^{K \times T \times S}$（因为 Adam 的 momentum/variance 是 per-parameter 的，天然支持不同 slice 的不同更新频率）。

\begin{itemize}[leftmargin=1.5em]
\item \textbf{Learning rate}：推荐 $\texttt{logits\_lr} \in [10^{-2}, 10^{-1}]$（logits 空间需要较大 lr）。
\item \textbf{Weight decay}：$0$（logits 不需要正则化，softmax 自带 bounded output）。
\item \textbf{Gradient clipping}：$\texttt{max\_grad\_norm} = 1.0$（防止单个高 advantage sample 导致 logits 跳变过大）。
\end{itemize}
\end{strategy}

%% ============================================================
\section{Reward 采样与评估}
\label{sec:reward_eval}
%% ============================================================

\subsection{Sampling 阶段}

\begin{algorithm}[H]
\caption{Per-Set MoTV Sampling}
\label{alg:sampling}
\begin{algorithmic}[1]
\Require EMA logits $\bar{\ell}$, dataloader (mixed prompt sets)
\Ensure Samples with rewards and prompt set labels
\For{each batch from dataloader}
    \For{each sample $b$ in batch}
        \State $s_b \gets \texttt{prompt\_set\_id}[b]$
        \State Generate $x_b$ using $\bar{\ell}_{:,:,s_b}$ (Algorithm~\ref{alg:inference})
    \EndFor
    \State Compute rewards: $R_{\mathrm{gen}}(x_b)$, $R_{\mathrm{pick}}(x_b)$, $R_{\mathrm{ocr}}(x_b)$ for all $b$
    \State Compute per-set reward: $\reward_{s_b}(x_b)$ using eq.~\eqref{eq:reward_def}
\EndFor
\end{algorithmic}
\end{algorithm}

\subsection{Reward Evaluation 开销}

\begin{remark}[三倍 Reward 计算]
每个 sample 需要计算\textbf{所有三个} reward function（不仅是自己 set 的），因为 OOD bonus 需要跨 set reward。对于仅 inference-based rewards（如 PickScore, GenEval），这是轻量级的 forward pass。OCR reward 需要 OCR model inference。总 reward 评估开销 $\approx 3\times$ 单 reward 成本。
\end{remark}

%% ============================================================
\section{理论对比与实验预期}
\label{sec:expected}
%% ============================================================

\subsection{对比总表}

\begin{center}
\renewcommand{\arraystretch}{1.5}
\footnotesize
\begin{tabular}{@{}p{2.5cm} c c c c c@{}}
\toprule
& \textbf{OPD-Avg} & \textbf{OPD-Route} & \textbf{MoTV (global)} & \textbf{MoTV (per-set)} \\
\midrule
Parameters & $\sim 10^7$ & $\sim 10^7$ & $K \times T$ & $K \times T \times S$ \\
In-domain 保证 & \ding{55} $< $ single & $\approx$ single & 弱（折衷） & \ding{51} $\geq$ single \\
Level~2 可能 & \ding{55} & \ding{55} & 中（受折衷限制） & \ding{51} 高 \\
Cross-set 干扰 & 强（共享 $\theta$） & 无（分离） & 有（共享 $\lambda$） & \textbf{无} \\
Reward needed & \ding{55} & \ding{55} & \ding{51} & \ding{51} \\
Optimal obj. & KL min & KL min & Reward max & Reward max \\
\bottomrule
\end{tabular}
\end{center}

\subsection{预期的 $\lambda$ 收敛模式}

\begin{tcolorbox}[colback=yellow!5, colframe=yellow!60!black, title=预期实验结果]
训练收敛后，我们预期看到以下 $\lambda$ pattern：

\textbf{Set = GenEval}（$\ell_{:,:,\mathrm{gen}}$）：
\begin{itemize}[leftmargin=1em]
\item 早期 timestep ($t_1$--$t_8$)：GenEval teacher 占主导（$\lambda_{\mathrm{gen}} \sim 0.8$--$0.9$），因为 layout/composition 在早期确定。
\item 晚期 timestep ($t_{20}$--$t_{28}$)：可能少量 PickScore teacher 贡献（$\lambda_{\mathrm{pick}} \sim 0.1$--$0.2$），因为清晰的 rendering 有助于 GenEval 的属性判断。
\end{itemize}

\textbf{Set = OCR}（$\ell_{:,:,\mathrm{ocr}}$）：
\begin{itemize}[leftmargin=1em]
\item 早期 timestep：可能 GenEval teacher 有贡献（文字 spatial positioning）。
\item 中期 timestep ($t_8$--$t_{18}$)：OCR teacher 占主导（笔画形成阶段）。
\item 晚期 timestep：OCR teacher 仍主导，可能少量 PickScore 提升非文字区域美感。
\end{itemize}

\textbf{Set = PickScore}（$\ell_{:,:,\mathrm{pick}}$）：
\begin{itemize}[leftmargin=1em]
\item 全程 PickScore teacher 占主导（$\lambda_{\mathrm{pick}} \sim 0.7$--$0.9$），因为 aesthetics 是全局属性。
\item 早期 timestep 可能少量 GenEval teacher 贡献（正确的 composition 本身提升美感）。
\end{itemize}
\end{tcolorbox}

%% ============================================================
\section{总结}
\label{sec:summary}
%% ============================================================

\begin{tcolorbox}[colback=green!5, colframe=green!50!black, title=Phase 2 核心设计总结]
\begin{enumerate}[leftmargin=1.5em]
\item \textbf{架构}：$\ell \in \R^{K \times T \times S}$（252 个参数），每个 prompt set 独立的 teacher 混合策略。

\item \textbf{Loss}：Per-set DiffusionNFT loss，各 set 完全解耦。梯度具有三重正交性（timestep × prompt-set × teacher-softmax）。

\item \textbf{Reward}：$r_{\mathrm{composite}}(x, s) = R_{\mathrm{in}(s)}(x) + \gamma \cdot \mathrm{mean}(R_{\mathrm{ood}})$，normalization 交给 GDPO-style advantage。

\item \textbf{Initialization}：Teacher-biased（$C=2$）从 near-single-teacher 出发，保证初始不退化。

\item \textbf{保证}：Level~1（$\geq$ single teacher）有构造性证明。Level~2（$>$ single teacher）由 OOD bonus 驱动探索。

\item \textbf{Expected outcome}：每个 prompt set 上至少匹配 single teacher in-domain performance，同时在 timestep 维度上发现 cross-teacher 互补性，实现 OOD 提升甚至 in-domain 超越。
\end{enumerate}
\end{tcolorbox}

%% ============================================================
\section{Reward-Advantage Pipeline: 理论推导与实现分析}
\label{sec:reward_advantage}
%% ============================================================

\subsection{Multi-Teacher / Multi-Reward / Multi-Dataset 对应关系}

\begin{center}
\renewcommand{\arraystretch}{1.4}
\begin{tabular}{@{}lccc@{}}
\toprule
\textbf{Source (dataset)} & \textbf{In-domain Teacher} & \textbf{In-domain Reward} & \textbf{Applicable Rewards} \\
\midrule
\texttt{geneval} & GenEval Teacher & \texttt{geneval} & geneval, pick\_score \\
\texttt{pickscore} & PickScore Teacher & \texttt{pick\_score} & pick\_score \\
\texttt{ocr} & OCR Teacher & \texttt{ocr} & ocr, pick\_score \\
\bottomrule
\end{tabular}
\end{center}

\begin{definition}[对应关系的三层结构]
\begin{enumerate}[leftmargin=1.5em]
\item \textbf{Teacher $\to$ Source}：由 \texttt{TeacherConfig.sources} 指定。Teacher $k$ 的 velocity 用于所有 source 的 sample（MoF 中所有 teacher 都 forward），但 \texttt{sources} 决定了哪个 teacher 是该 source 的``in-domain teacher''（用于 \texttt{teacher\_biased} 初始化）。

\item \textbf{Source $\to$ In-domain Reward}：由 \texttt{TeacherConfig.reward\_name} 显式指定。决定了对 source $s$ 的 sample 计算 composite reward 时，哪个 reward 是主信号（权重 = 1）。

\item \textbf{Reward $\to$ Applicable Sources}：由 \texttt{RewardConfig.applicable\_sources} 指定。决定了哪些 source 的 sample 可以被该 reward model 评分（不适用的 sample 得到 NaN）。
\end{enumerate}
\end{definition}

\subsection{Reward Aggregation 公式（旧方案 vs.\ 正确方案）}

\paragraph{旧方案（raw reward 空间聚合 → 再 normalize）：}
\begin{equation}
r_{\mathrm{composite}}(x, s) = R_{\mathrm{in}(s)}(x) + \gamma \cdot \mathrm{mean}(R_{\mathrm{ood}})
\quad\xrightarrow{\text{group norm}}\quad
\mathcal{A}_i = \frac{r_i - \mu_g}{\sigma_g \cdot \sigma_{\mathrm{global}}}
\tag{旧}
\end{equation}
\textbf{问题}：GenEval $\in [0,1]$ 与 PickScore $\in [0.8, 0.95]$ 的 scale 差异使 $\gamma$ 的有效强度不可控。

\paragraph{正确方案（per-reward normalize → advantage 空间聚合）：}

与 Section~4.2 的 Definition~\ref{def:advantage} 一致：
\begin{align}
&\text{Step 1:}\quad a_k(x_i) = \frac{R_k(x_i) - \mu_{k,g}}{\sigma_{k,g}} \quad\text{(per-reward, per-group)} \tag{eq.~\ref{eq:per_reward_norm}} \\
&\text{Step 2:}\quad \mathcal{A}_{\mathrm{combined}}(x_i) = a_{\mathrm{in}(s)}(x_i) + \gamma \cdot \mathrm{mean}(a_{\mathrm{ood}}) \tag{eq.~\ref{eq:advantage_aggregation}} \\
&\text{Step 3:}\quad \mathcal{A}_{\mathrm{final}}(x_i) = \mathcal{A}_{\mathrm{combined}}(x_i) / \sigma_{\mathrm{global}} \tag{eq.~\ref{eq:final_norm}}
\end{align}

此方案中 $\gamma = 0.2$ 的语义清晰：OOD advantage（已归一化为 $\sim N(0,1)$）贡献 20\% 的信号强度，与 reward 原始 scale 完全无关。

\subsection{具体示例（3-teacher config）}

\begin{itemize}[leftmargin=1.5em]
\item \texttt{geneval} sample: $\mathcal{A} = a_{\mathrm{geneval}} + 0.2 \cdot a_{\mathrm{pick\_score}}$
\item \texttt{pickscore} sample: $\mathcal{A} = a_{\mathrm{pick\_score}}$（无 OOD applicable）
\item \texttt{ocr} sample: $\mathcal{A} = a_{\mathrm{ocr}} + 0.2 \cdot a_{\mathrm{pick\_score}}$
\end{itemize}

其中每个 $a_k$ 已经是 group-normalized（zero-mean, unit-variance within group）。

\subsection{分布式训练下的正确性要求}

\begin{tcolorbox}[colback=red!5, colframe=red!50, title=关键约束]
在 \texttt{distributed\_k\_repeat} sampler 下，同一 group 的 $K$ 个 sample 可能分散在不同 rank 上。正确的 group-mean 需要\textbf{跨 rank gather}——只用本地 samples 计算 group-mean 会得到错误结果。

当前实现通过 \texttt{AdvantageProcessor.collect\_group\_rewards()} 解决：
\begin{enumerate}[leftmargin=1.5em]
\item Pack rewards + unique\_id → \texttt{(B, 2)} tensor
\item \texttt{accelerator.gather()} → \texttt{(W*B, 2)} global tensor
\item 在 global tensor 上计算 group-mean / global-std
\item \texttt{\_to\_local()} scatter 回各 rank
\end{enumerate}
\end{tcolorbox}

\subsection{实现中的潜在问题分析}

\begin{enumerate}[leftmargin=1.5em]

\item \textbf{Cross-source Group 混合问题}（中等风险）

当前 sampling 使用 \texttt{\_interleaved\_source\_iter()} 交替产出不同 source 的 batch。每个 batch 对应一个 source，batch 内的 samples 共享相同的 \texttt{unique\_id}（同 prompt 的 $K$ 个 sample 在同一 batch 内）。

\textbf{但}：如果 \texttt{group\_size=16} 且 \texttt{per\_device\_batch\_size=8}，则一个 group 需要 2 个 batch 来生成 16 个 sample（batch size $\times$ 8 GPU $= 64$ per batch，但 unique\_sample\_num\_per\_epoch 控制不同 prompt 的数量）。

\textbf{关键问题}：同一 group 内的所有 sample 必须来自同一 source（因为它们共享同一个 prompt）。当前的 \texttt{\_interleaved\_source\_iter()} 保证每个 batch 是单 source 的，但同一 prompt 的 $K$ 个重复 sample 可能跨多个 batch。由于 interleaved 模式 round-robin 切换 source，这些重复 sample 确实在同一个 source 的 batch 中生成——\textbf{正确，无问题}。

\item \textbf{Group-Std vs.\ Global-Std 的语义差异}（低风险但需注意）

当前实现做了两层 normalization：先 group-std（eq.~\ref{eq:group_norm}），再 global-std（eq.~\ref{eq:global_norm}）。这等价于：
\begin{equation}
\mathcal{A}_i = \frac{r_i - \mu_g}{\sigma_g \cdot \sigma_{\mathrm{global}}}
\end{equation}

如果不同 source 的 group 内 reward variance 差异很大（如 geneval 的 group-std $\approx 0.3$，pickscore 的 group-std $\approx 0.02$），group-std normalization 后它们的 scale 已经对齐。global-std 此时主要作用是稳定整体 scale（不改变相对大小）。

\textbf{替代方案}：只用 group-mean 减去（不除 group-std），然后用 global-std 统一 normalize。这保留了不同 source 的 reward 信号强度差异，可能更好：
\begin{equation}
\mathcal{A}_i^{\mathrm{alt}} = \frac{r_i - \mu_g}{\sigma_{\mathrm{global}}}
\end{equation}

但当前的双层 normalization 在 GDPO 论文中已验证有效，暂不修改。

\item \textbf{Composite Reward 的 Scale 不一致}（需注意）

由于 \texttt{applicable\_sources} 过滤，不同 source 的 sample 的 composite reward 组成不同：
\begin{itemize}
\item geneval sample: $R_{\mathrm{geneval}} + 0.2 \cdot R_{\mathrm{pick}}$（两项之和）
\item pickscore sample: $R_{\mathrm{pick}}$（单项）
\end{itemize}

两者的原始 scale 不同。但由于 group-mean 是\textbf{组内}计算（同一 prompt $\Rightarrow$ 同一 source $\Rightarrow$ 同一 composite 公式），且 group-std 归一化后 scale 对齐，这\textbf{不是问题}。

\item \textbf{\texttt{\_to\_local()} 返回类型}（minor bug）

\texttt{AdvantageProcessor.\_to\_local()} 在 \texttt{group\_on\_same\_rank=True} 时返回 \texttt{torch.Tensor}（可能在 GPU 上），在 \texttt{group\_on\_same\_rank=False} 时返回 CPU numpy array 经 \texttt{torch.as\_tensor} 转换。后续 \texttt{float(adv)} 对 GPU tensor 有隐式 D2H transfer。建议统一 \texttt{.cpu().numpy()} 后再 store。

\end{enumerate}

\subsection{推荐的实现改进}

\begin{tcolorbox}[colback=green!5, colframe=green!50!black, title=短期可执行的改进]
\begin{enumerate}[leftmargin=1.5em]
\item \textbf{fix}: 在 \texttt{prepare\_feedback} 中 \texttt{local\_advantages} 转为 numpy 后再逐 sample 存储（避免逐个 GPU→CPU transfer）。

\item \textbf{enhancement}: 为 pickscore source 也配置 OOD reward（如 geneval），使其也能受益于 cross-teacher signal：
\begin{verbatim}
- name: "geneval"
  applicable_sources: [geneval, pickscore]
\end{verbatim}
这样 pickscore sample 的 composite = $R_{\mathrm{pick}} + 0.2 \cdot R_{\mathrm{geneval}}$，获得更丰富的梯度信号。

\item \textbf{monitoring}: 添加 per-source advantage mean 的 logging，确认各 source 的 advantage 分布是否平衡（不应有某个 source 的 advantage 系统性偏高或偏低）。
\end{enumerate}
\end{tcolorbox}

\end{document}
